% === Preface ===
\section*{Preface}
\addcontentsline{toc}{section}{Preface}

\subsection*{Why This Book}

Post-training is the most rapidly evolving area in current research on large language models (LLMs) and vision-language models (VLMs).
From the establishment of the SFT$\rightarrow$RLHF pipeline by InstructGPT in 2022,
to DeepSeek-R1 demonstrating in 2025 that pure RL can give rise to emergent reasoning capabilities from scratch,
to thinking models and computer use becoming the default industry paradigm in 2026---
in just four years, the core methodology of post-training has undergone multiple fundamental paradigm shifts.

Existing surveys predominantly adopt a static taxonomy, organizing content by method type (SFT, RLHF, DPO) or capability dimension (reasoning, safety, multimodal).
Such an organizational scheme is clear at any single temporal cross-section, but struggles to answer a more fundamental question:
\textbf{Why did these methods emerge? What problems did they solve? And what new problems did they give rise to?}
The motivation for writing this book stems precisely from this observation---
we aim to re-examine the post-training landscape from an \textbf{evolutionary} perspective,
tracing the developmental chain of ``problem $\rightarrow$ solution $\rightarrow$ new problem $\rightarrow$ improved solution,''
and revealing the flow of ideas across different research directions.

\subsection*{How to Read This Book}

This book organizes the post-training landscape into seven problem-driven evolutionary threads,
each corresponding to a core research question. Below are several recommended reading paths:

\begin{itemize}[leftmargin=2em]
  \item \textbf{Linear reading} (Thread 1 $\rightarrow$ Thread 7):
    Read all chapters sequentially for a complete panorama of the post-training field.
    The chapters follow a logical progression---SFT lays the foundation, preference optimization supplements preference signals,
    safety imposes constraints, reasoning extends capability boundaries, multimodal broadens modalities,
    agentic connects to the real world, and efficiency brings everything to deployment.

  \item \textbf{Thread-based reading}:
    Select any thread of interest and read it directly,
    using the colored cross-reference arrows in the text ({\color{thread2}$\rightarrow$} \textit{see \S X.Y} format) to jump to corresponding sections in related threads.
    Each thread is self-contained, but the cross-references reveal deeper connections.

  \item \textbf{Global navigation}:
    Figure~1 in the Introduction provides a panoramic evolution map of all seven threads,
    serving as a navigation guide---first locate the region of interest, then dive into the corresponding chapter.
\end{itemize}

\noindent
This book employs a unified box system to highlight key information:
\begin{itemize}[leftmargin=2em]
  \item \textbf{\color{red!60}Red boxes} (Problem Box): Mark core problems that drive the evolution of the field
  \item \textbf{\color{blue!60}Blue boxes} (Solution Box): Present key solutions to the identified problems
  \item \textbf{\color{green!60!black}Green boxes} (Evolution Box): Trace evolutionary paths of methodologies
  \item \textbf{\color{yellow!70!black}Yellow boxes} (Open Problem Box): Highlight unresolved open problems
\end{itemize}

\subsection*{Limitations of This Book}

This book has the following known limitations, which readers should bear in mind:

\begin{enumerate}[leftmargin=2em]
  \item \textbf{Knowledge cutoff}: The literature coverage of this book extends through early 2026.
    Given the extremely rapid pace of evolution in the post-training field, some discussions may require updating by the time of publication.

  \item \textbf{Method-oriented focus}: This book focuses on the evolution of training methodologies,
    with relatively limited coverage of engineering deployment (serving, inference optimization) and productization practices.

  \item \textbf{Literature coverage}: This book primarily draws on published academic work and publicly available technical reports.
    A substantial body of industrial practice remains not fully covered,
    which means that the real-world application scenarios for certain methods may be broader than described herein.
\end{enumerate}

\subsection*{Acknowledgments}

We are grateful to all researchers who have contributed to the post-training field---
it is their work that has made the evolution of this area so rich and compelling.
Throughout the writing of this book, discussions and exchanges with colleagues have provided many valuable perspectives and insights.
